%=============================================================================
\chapter{Multivariate heteroskedastic BART}
\label{ch:full}

This chapter derives, from scratch, every step of the BART sampler for the
most general model treated in this document:
\begin{itemize}
	\item each outcome $\vy_i$ is a $k$-vector;
	\item each leaf parameter $\vmu_\ell$ is a $k$-vector, with a centred
	      multivariate-Gaussian prior of precision $\mLamm$;
	\item the noise on $\vy_i$ has precision $\mA_i\,\mLam\,\mA_i$ where
	      $\mA_i = \diag\va_i$ is a fixed (known) diagonal matrix of
	      nonnegative datapoint-specific precision scales, and $\mLam$ is a
	      global $k \times k$ error precision;
	\item $\mLam$ has a Wishart conjugate prior;
	\item trees follow the Chipman--George--McCulloch (CGM) prior and are
	      updated by GROW/PRUNE moves only; there is no variable-selection
	      prior.
\end{itemize}
Working in the multivariate setting it is more convenient to deal with
precision matrices than with covariance matrices, so throughout this chapter
$\mLam$ and $\mLamm$ are the primary parameters: all Gaussian distributions
are written in precision form $\N_k(\cdot,\, \mLam^{-1})$ and the conjugate
prior on $\mLam$ is a Wishart (rather than an inverse-Wishart on the
covariance $\mLam^{-1}$). Using precision scales $\va_i$ in place of weights
also relaxes positivity: $a_{i,j} = 0$ is admissible and corresponds to a
missing $j$-th component of $\vy_i$ (infinite noise on that coordinate).
Following bartz, when some components can be missing, $\mLam$ is restricted
to be diagonal (\S\ref{sec:mv-model}).

%-----------------------------------------------------------------------------
\section{Notation and data}
\label{sec:mv-notation}

We observe $n$ pairs $(\vx_i, \vy_i)_{i=1}^n$ with $\vx_i \in \R^p$ and
$\vy_i \in \R^k$. Stack the responses as rows into the matrix
$Y \in \R^{n \times k}$ and the covariates row-by-row into the design matrix
$X = (\vx_1, \dots, \vx_n)^\top \in \R^{n \times p}$. In addition, each
datapoint comes with a vector of nonnegative precision scales
$\va_i \in \R^k_{\geq 0}$, treated as fixed and known (part of the
observation); we collect them in $\mA_i = \diag\va_i$.

A \emph{decision tree} $T$ is a binary tree whose internal nodes carry a
splitting rule of the form ``$x_{v} < c$'', for some variable index
$v \in \{1, \dots, p\}$ and cutpoint $c$. Given $T$ and a covariate $\vx$, the
tree assigns $\vx$ to a unique leaf $\ell(\vx; T)$ by following the rules from
the root. We write $L(T)$ for the set of leaves of $T$ and
$|L(T)|$ for their number. To each tree we attach a vector-valued leaf array
\begin{equation}
	\label{eq:mv-leaf-array}
	M = \{\, \vmu_\ell : \ell \in L(T) \,\} \subset \R^k,
\end{equation}
and define the corresponding $k$-dimensional regression function
\begin{equation}
	\label{eq:mv-g}
	g(\vx; T, M) = \sum_{\ell \in L(T)} \vmu_\ell\, \Ind[\ell(\vx; T) = \ell]
	\in \R^k.
\end{equation}
A BART model uses a sum of $m$ such trees,
\begin{equation}
	\label{eq:mv-sum-of-trees}
	f(\vx) = \sum_{t=1}^m g(\vx;\, T_t, M_t) \in \R^k.
\end{equation}

%-----------------------------------------------------------------------------
\section{The generative model}
\label{sec:mv-model}

\subsection{Conditional outcome distribution}

Conditional on the trees and the noise precision, the outcomes are
\begin{equation}
	\label{eq:mv-lik}
	\vy_i \mid T_{1:m}, M_{1:m}, \mLam \stackrel{\text{ind}}{\sim}
	\N_k\!\left(f(\vx_i),\; (\mA_i\,\mLam\,\mA_i)^{-1}\right),
	\qquad i = 1,\dots,n.
\end{equation}
Equivalently, in terms of rescaled residuals,
$\va_i \odot (\vy_i - f(\vx_i)) \sim \N_k(\mathbf{0}, \mLam^{-1})$. We will
use this form repeatedly.

Viewed as a function of the parameters $(T_{1:m}, M_{1:m}, \mLam)$ with
$Y$ fixed, the product over $i$ is the \emph{likelihood} of the model. We
will use the word ``likelihood'' only in this strict sense, and only in
contexts where the parameters (or some marginal thereof) are the free
variables.

\paragraph{Missing components.}

A component $y_{i,j}$ with $a_{i,j} = 0$ has infinite noise, i.e., it is
missing: the information carried by $\vy_i$ is the marginal distribution of
its other components. Let
\begin{equation}
	\label{eq:mv-Oi}
	O_i \coloneqq \{j : a_{i,j} > 0\},
	\qquad
	k_i \coloneqq |O_i|,
\end{equation}
be the observed components of $\vy_i$ and their number, and $\bar O_i$ the
missing ones. Marginalising \eqref{eq:mv-lik} over $\vy_{i,\bar O_i}$, the
covariance of $\vy_{i,O_i}$ is the $O_i$ block of
$\mA_i^{-1}\mLam^{-1}\mA_i^{-1}$, so
\begin{equation}
	\label{eq:mv-lik-obs}
	\begin{gathered}
		\vy_{i,O_i} \mid T_{1:m}, M_{1:m}, \mLam \stackrel{\text{ind}}{\sim}
		\N_{k_i}\left(f(\vx_i)_{O_i},
		\left(\mA_{i,O_i}\left[(\mLam^{-1})_{O_iO_i}\right]^{-1}\mA_{i,O_i}\right)^{-1}\right),
		\\
		\mA_{i,O_i} \coloneqq \diag\va_{i,O_i},
	\end{gathered}
\end{equation}
where $[(\mLam^{-1})_{O_iO_i}]^{-1} = \mLam_{O_iO_i} -
	\mLam_{O_i\bar O_i}\mLam_{\bar O_i\bar O_i}^{-1}\mLam_{\bar O_iO_i}$ is the
Schur complement of $\mLam_{\bar O_i\bar O_i}$ in $\mLam$. This does not
depend on the missing scales, so it is also the limit $a_{i,j} \to 0$,
$j \in \bar O_i$, of the complete model. The precision in
\eqref{eq:mv-lik-obs} depends on $\mLam$ and on the missingness pattern of
$i$ in a way that does not factor as $\mLam$ times a data-only matrix, which
would break the sufficient statistics of \S\ref{sec:mv-partial-resid} and the
conjugacy of \S\ref{sec:mv-Sigma}. For this reason, following bartz,
\emph{whenever some component can be missing we restrict $\mLam$ to be
	diagonal}. Then $[(\mLam^{-1})_{O_iO_i}]^{-1} = \mLam_{O_iO_i}$, and
\eqref{eq:mv-lik-obs} becomes the Gaussian whose precision is
$\mA_i\mLam\mA_i$ restricted to the observed rows and columns,
\begin{equation}
	\label{eq:mv-lik-diag}
	\begin{gathered}
		\vy_{i,O_i} \mid T_{1:m}, M_{1:m}, \mLam \stackrel{\text{ind}}{\sim}
		\N_{k_i}\left(f(\vx_i)_{O_i},
		[\mA_i \mLam \mA_i]_{O_iO_i}^{-1}\right),
		\\
		[\va_i \odot (\vy_i - f(\vx_i))]_{O_i} \sim
		\N_{k_i}(\mathbf{0}, \mLam_{O_iO_i}^{-1}).
	\end{gathered}
\end{equation}
The missing entries of $\vy_i$ may be set to any value, since they are
always multiplied by $a_{i,j} = 0$. In what follows, formulas written with
the full $k \times k$ precision $\mA_i\mLam\mA_i$, whose missing rows and
columns vanish, are to be read as \eqref{eq:mv-lik-diag}.

Many formulas written without taking into account that $a$ can be 0 continue to
be correct in the missing data case, with the null components of $\mathbf a_i$
implicitly deleting the terms that ought to disappear (again, only under the
restriction that $\mLam$ be diagonal). The implementation takes advantage of
this to simplify the code. In some places this does not hold, and we will point
it out.

\subsection{Tree-structure prior}
\label{sec:mv-tree-prior}

A priori, each tree $T_t$ is drawn independently from the CGM process. Working
node-by-node from the root, a node at depth $d \ge 0$ (with the root at depth
$0$) is, conditionally on existing, internal (i.e., nonterminal) with
probability
\begin{equation}
	\label{eq:pnt}
	p_{\mathrm{NT}}(d) = \alpha\,(1 + d)^{-\beta},
	\qquad \alpha \in (0,1),\ \beta > 0,
\end{equation}
and otherwise a leaf. The implementation truncates the recursion at a maximum
depth $d_{\max}$, enforced by setting $p_{\mathrm{NT}}(d_{\max}) = 0$ (in
\texttt{bartz/mcmcstep/\_state.py}, \texttt{make\_p\_nonterminal} evaluates
\eqref{eq:pnt} at depths $0, \dots, d_{\max} - 1$ and \texttt{init} appends
the final $0$ through \texttt{parse\_p\_nonterminal}).

Conditionally on a node being internal, its splitting rule is drawn from
\begin{itemize}
	\item a uniform distribution on the variables $v$ that still admit at
	      least one valid cutpoint at that node, and
	\item a uniform distribution on the cutpoints $c$ that are compatible with
	      the ancestor splits along the path to the root.
\end{itemize}
Each predictor $v$ comes with a fixed list of allowed cutpoints, specified as
part of the model (separate from, though often derived by default from, the
values $X_{\cdot,v}$ observed in the training set). ``Compatible'' means
that the cutpoint lies strictly between the inherited lower and upper bounds
imposed by the ancestor splits on variable $v$ along the path to the root.
In \texttt{bartz} the allowed cutpoints are pre-binned and live in the
integer range $[1,\, \texttt{max\_split}_v]$ (see the
\texttt{bartz.prepcovars} module);
the prior on the rule is uniform on the available $(v,c)$ pairs.

We write $p(T_t)$ for the marginal probability of the tree structure
(including the splitting rules) under this prior.

\subsection{Leaf prior}

Conditional on the tree, the leaf parameters are i.i.d.\ centred multivariate
Gaussian in precision form:
\begin{equation}
	\label{eq:mv-leaf-prior}
	\vmu_\ell \mid T_t \stackrel{\text{i.i.d.}}{\sim}
	\N_k\!\left(\mathbf{0},\, \mLamm^{-1}\right),
	\qquad \ell \in L(T_t).
\end{equation}
The precision $\mLamm$ is calibrated so that the prior on the sum of trees,
$f(\vx) \sim \N_k(\mathbf{0}, m\,\mLamm^{-1})$, has user-controlled
variability.

\subsection{Error precision prior}

The global noise precision follows the conjugate Wishart prior
\begin{equation}
	\label{eq:mv-iw-prior}
	\mLam \sim \W(\nu,\, \Psi),
\end{equation}
parametrised so that $\Psi \in \R^{k \times k}$ is the prior inverse-scale
(rate) matrix and $\nu > k - 1$ the degrees of freedom; the prior density is
proportional to $|\mLam|^{(\nu - k - 1)/2}\exp\!\left(-\frac{1}{2}
	\tr(\Psi\,\mLam)\right)$.

\subsection{Bayesian network}
\label{sec:mv-bn}

Putting the pieces together, the joint density factorises as
\begin{equation}
	\label{eq:mv-joint}
	p(Y, T_{1:m}, M_{1:m}, \mLam)
	= p(\mLam) \prod_{t=1}^m p(T_t)\, p(M_t \mid T_t)
	\cdot \prod_{i=1}^n p(\vy_i \mid T_{1:m}, M_{1:m}, \mLam),
\end{equation}
with $p(M_t \mid T_t) = \prod_{\ell \in L(T_t)} \N_k(\vmu_\ell;\,
	\mathbf{0}, \mLamm^{-1})$.

Graphically:
\begin{center}
	\begin{tikzpicture}[node distance=0.9cm and 1.6cm,
			every node/.style={font=\footnotesize}]
		% Top row: hyperparameters
		\node[const]                       (alpha) {$\alpha,\beta$};
		\node[const, right=2.2cm of alpha] (sm)    {$\mLamm$};
		\node[const, right=2.2cm of sm]    (nu)    {$\nu,\Psi$};
		% Middle row: latent variables
		\node[latent, below=of alpha] (T)  {$T_t$};
		\node[latent, below=of sm]    (mu) {$\vmu_\ell^{(t)}$};
		\node[latent, below=of nu]    (S)  {$\mLam$};
		% Bottom row: data
		\node[obs, below=1.6cm of mu] (y) {$\vy_i$};
		\node[const, left=of y]       (x) {$\vx_i$};
		\node[const, right=of y]      (w) {$\va_i$};

		\edge {alpha} {T};
		\edge {sm}    {mu};
		\edge {nu}    {S};
		\edge {T}     {mu};
		\edge {mu}    {y};
		\edge {S}     {y};
		\edge {x}     {y};
		\edge {w}     {y};
		\draw[->, >={triangle 45}] (T) to[bend right=25] (y);

		\plate {leaves} {(mu)}        {$\ell \in L(T_t)$};
		\plate {trees}  {(T)(leaves)} {$t = 1,\dots,m$};
		\plate {data}   {(x)(y)(w)}   {$i = 1,\dots,n$};
	\end{tikzpicture}
\end{center}
The dependence of $\vy_i$ on the tree variables is mediated by the assignment
$\ell(\vx_i; T_t)$, which is a deterministic function of $\vx_i$ and $T_t$.
The precision scales $\va_i$ enter as fixed, observed inputs alongside
$\vx_i$.

%-----------------------------------------------------------------------------
\section{Single-tree decomposition and partial residuals}
\label{sec:mv-partial-resid}

Bayesian backfitting handles \eqref{eq:mv-sum-of-trees} by updating one tree
at a time, holding the others fixed. Define the \emph{partial residual} for
tree $t$ by removing every other tree's contribution:
\begin{equation}
	\label{eq:mv-partial-resid}
	\vr_i^{(t)} \coloneqq \vy_i - \sum_{k \ne t} g(\vx_i;\, T_k, M_k)
	\in \R^k.
\end{equation}
By \eqref{eq:mv-lik},
\begin{equation}
	\label{eq:mv-partial-lik}
	\vr_i^{(t)} \mid T_t, M_t, \mLam \stackrel{\text{ind}}{\sim}
	\N_k\!\left(g(\vx_i;\, T_t, M_t),\; (\mA_i\,\mLam\,\mA_i)^{-1}\right),
\end{equation}
restricted to the observed components as in \eqref{eq:mv-lik-diag}, so the
conditional update of $(T_t, M_t)$ is identical to a single-tree
regression on the partial-residual data $\{(\vx_i, \va_i, \vr_i^{(t)})\}_{i=1}^n$.

For the rest of the chapter we fix $t$ and drop the superscript:
$\vr_i = \vr_i^{(t)}$, $T = T_t$, $M = M_t$. For each leaf $\ell \in L(T)$
define the index set and count
\begin{equation}
	\label{eq:mv-Iell}
	I_\ell = \{\, i : \ell(\vx_i; T) = \ell \,\},
	\qquad
	n_\ell = |I_\ell|.
\end{equation}
By construction $\sum_{\ell} n_\ell = n$ and the index sets $\{I_\ell\}$
partition $\{1,\dots,n\}$.

The per-datapoint precision matrix factors as a Hadamard product with an
outer product. Define
\begin{equation}
	\label{eq:mv-Atilde}
	\tilde\mA_i \;\coloneqq\; \va_i\va_i^\top \;\in\;\R^{k\times k};
\end{equation}
then
\begin{equation}
	\label{eq:mv-Pi}
	\mP_i \;\coloneqq\; \mA_i\,\mLam\,\mA_i
	= \mLam \odot \tilde\mA_i
	\;\in\;\R^{k\times k},
\end{equation}
with $\odot$ the Hadamard (entrywise) product, since
$[\mA_i\mLam\mA_i]_{ab} = a_{i,a}\Lambda_{ab}a_{i,b}
	= \Lambda_{ab}\,[\tilde\mA_i]_{ab}$.

Because the per-datapoint precision varies with $i$, the per-leaf sufficient
statistics are not scalars but $k\times k$ matrices. Define the
$\mLam$-dependent posterior precision contribution and score by
\begin{align}
	\label{eq:mv-precision-sum}
	\mP_\ell & \coloneqq \sum_{i \in I_\ell}\mP_i
	= \sum_{i \in I_\ell}\mLam \odot \tilde\mA_i
	= \mLam \odot \mN_\ell
	\;\in\;\R^{k\times k},
	\\
	\label{eq:mv-suffstats-B}
	\vB_\ell & \coloneqq \sum_{i \in I_\ell}\mP_i\,\vr_i
	= (\mLam \odot S_\ell)\,\Ind_k
	\;\in\;\R^{k},
\end{align}
where $\Ind_k$ is the all-ones $k$-vector and the data-only sufficient
statistics are
\begin{align}
	\label{eq:mv-suffstats-N}
	\mN_\ell & \coloneqq \sum_{i \in I_\ell}\tilde\mA_i
	\;\in\;\R^{k\times k},
	\\
	\label{eq:mv-suffstats-S}
	S_\ell   & \coloneqq \sum_{i \in I_\ell}\tilde\mA_i \odot (\Ind_k\,\vr_i^\top)
	\;\in\;\R^{k\times k}.
\end{align}
Both $\mN_\ell, S_\ell$ depend only on the data (not on $\mLam$) and
accumulate additively across any partition of the index set. $\mN_\ell$ is
symmetric; $S_\ell$ in general is not. $\mP_\ell$ and $\vB_\ell$ depend on
$\mLam$ only and are computable from the $O(k^2)$ data-only statistics
$(\mN_\ell, S_\ell)$.

In the homoskedastic special case $\va_i \equiv \Ind_k$,
$\mN_\ell = n_\ell\,\Ind_k\Ind_k^\top$, the $a$-th row of $S_\ell$ equals
$\sum_{i \in I_\ell}\vr_i^\top$, $\mP_\ell = n_\ell\,\mLam$, and
$\vB_\ell = \mLam\sum_{i \in I_\ell}\vr_i$.

%-----------------------------------------------------------------------------
\section{Full conditional of the leaves}
\label{sec:mv-full-cond}

Conditional on $T$ and $\mLam$, the $n_\ell$ residuals in leaf $\ell$ are
independent Gaussian,
$\vr_{i,O_i} \sim \N_{k_i}(\vmu_{\ell,O_i}, [\mP_i]_{O_iO_i}^{-1})$ by
\eqref{eq:mv-lik-diag}; since the missing rows and columns of $\mP_i$
vanish, the log density is
$-\frac{1}{2}(\vr_i - \vmu_\ell)^\top\mP_i(\vr_i - \vmu_\ell)$ up to a
constant, as if $\vr_i$ were a complete $k$-vector with mean $\vmu_\ell$ and
precision $\mP_i$. Across leaves the residuals are independent because
$\{I_\ell\}$ partitions the data. Combined with $\vmu_\ell \sim \N_k(\mathbf{0}, \mLamm^{-1})$, this is a
Gaussian--Gaussian conjugate update with heterogeneous data precisions.

\begin{prop}[Leaf full conditional]
	\label{prop:mv-leaf}
	For each $\ell \in L(T)$,
	\begin{equation}
		\label{eq:mv-leaf-fc}
		\vmu_\ell \mid T, \mLam, \vr_{1:n} \sim
		\N_k\!\left(
		\underbrace{\left(\mLamm + \mP_\ell\right)^{-1} \vB_\ell}_{=\,\hat\vmu_\ell},
		\;
		\underbrace{\left(\mLamm + \mP_\ell\right)^{-1}}_{=\,V_\ell}
		\right),
	\end{equation}
	independently across $\ell$.
\end{prop}

\begin{proof}
	The log of the joint of the data in leaf $\ell$ and the leaf prior, as a
	function of $\vmu_\ell$, is
	\begin{align}
		 & \sum_{i\in I_\ell}\log p(\vr_i \mid \vmu_\ell, \mLam)
		+ \log p(\vmu_\ell \mid \mLamm)                          \\
		 & = -\frac{1}{2}\sum_{i\in I_\ell}
		(\vr_i - \vmu_\ell)^\top \mP_i (\vr_i - \vmu_\ell)
		-\frac{1}{2}\vmu_\ell^\top \mLamm\, \vmu_\ell + \text{const.}
	\end{align}
	Expanding the quadratic form in $\vmu_\ell$,
	\begin{equation}
		(\vr_i - \vmu_\ell)^\top \mP_i (\vr_i - \vmu_\ell)
		= \vr_i^\top \mP_i \vr_i - 2\,\vmu_\ell^\top \mP_i \vr_i
		+ \vmu_\ell^\top \mP_i \vmu_\ell,
	\end{equation}
	and summing over $i \in I_\ell$, the $\vmu_\ell$-dependent terms collect,
	using \eqref{eq:mv-precision-sum}--\eqref{eq:mv-suffstats-B}, as
	\begin{equation}
		-\frac{1}{2}\vmu_\ell^\top
		\left(\mLamm + \mP_\ell\right) \vmu_\ell
		+ \vmu_\ell^\top \vB_\ell.
	\end{equation}
	Completing the square with $V_\ell^{-1} = \mLamm + \mP_\ell$ gives
	\begin{equation}
		-\frac{1}{2}(\vmu_\ell - \hat\vmu_\ell)^\top V_\ell^{-1}
		(\vmu_\ell - \hat\vmu_\ell) + \text{const},
		\qquad \hat\vmu_\ell = V_\ell \vB_\ell,
	\end{equation}
	yielding \eqref{eq:mv-leaf-fc}. Independence across leaves is immediate
	from the partition structure of $\{I_\ell\}$ and the independent leaf
	priors.
\end{proof}

In bartz, \eqref{eq:mv-leaf-fc} is computed by
\texttt{\_precompute\_leaf\_terms\_mv\_het} in \texttt{mcmcstep/\_step.py}:
$\mN_\ell$ is accumulated in \texttt{Forest.prec\_tree} as the per-leaf sum
of \texttt{State.prec\_scale} $= \tilde\mA_i$, the Cholesky factor of
$\mLamm + \mLam \odot \mN_\ell$ is stored in \texttt{mean\_factor}, and
$\vB_\ell$ is formed by \texttt{compute\_B} from the per-leaf sum of
$\tilde\mA_i \odot (\Ind_k\vr_i^\top)$ returned by \texttt{sum\_resid}. When
$\va_i \equiv \Ind_k$ (or $\va_i = a_i\Ind_k$ with a scalar scale per
datapoint), $\mP_\ell = n_\ell \mLam$ (resp.\ $\sum_{i\in I_\ell} a_i^2 \mLam$)
and the cheaper \texttt{\_precompute\_leaf\_terms\_mv} is used instead, with
the count from \texttt{Forest.count\_tree}.

%-----------------------------------------------------------------------------
\section{Marginal density with leaves integrated out}
\label{sec:mv-marg-lik}

Because the leaves are conjugate, we can integrate $M$ out analytically to
obtain a collapsed target on tree structure alone. Stack the residuals
belonging to leaf $\ell$ row-by-row into the matrix
\begin{equation}
	\label{eq:mv-Rell}
	\mR_\ell \coloneqq \{\,\vr_i^\top : i \in I_\ell\,\} \in \R^{n_\ell \times k}.
\end{equation}
Conditional on $T$ and $\mLam$, the blocks $\mR_\ell$ are independent across
$\ell$ (the leaves partition the data), so the marginal density of the whole
residual set factorises over leaves.

\begin{prop}[Per-leaf marginal density]
	\label{prop:mv-marglik}
	For each leaf $\ell$,
	\begin{align}
		\log p(\mR_\ell \mid T, \mLam)
		 & = -\frac{\log(2\pi)}{2}\sum_{i \in I_\ell} k_i
		+ \frac{1}{2}\sum_{i \in I_\ell}\log\left|[\mA_i\mLam\mA_i]_{O_iO_i}\right|
		\label{eq:mv-marglik-a}            \\
		 & \quad
		+ \frac{1}{2}\log|\mLamm|
		- \frac{1}{2}\log\bigl|\mLamm + \mP_\ell\bigr|
		\label{eq:mv-marglik-b}            \\
		 & \quad
		- \frac{1}{2}\sum_{i \in I_\ell} \vr_i^\top \mP_i \vr_i
		+ \frac{1}{2}\,\vB_\ell^\top \bigl(\mLamm + \mP_\ell\bigr)^{-1}
		\vB_\ell,
		\label{eq:mv-marglik-c}
	\end{align}
	and the joint log marginal of the tree is
	$\log p(\mR_{1:|L(T)|} \mid T, \mLam) = \sum_{\ell \in L(T)}
		\log p(\mR_\ell \mid T, \mLam)$.
\end{prop}

\begin{proof}
	By marginalisation over $\vmu_\ell$,
	\begin{equation}
		p(\mR_\ell \mid T, \mLam)
		= \int_{\R^k} p(\mR_\ell \mid \vmu_\ell, \mLam)\,
		p(\vmu_\ell \mid \mLamm)\, d\vmu_\ell.
	\end{equation}
	Substituting the explicit Gaussian densities (the data factor in the
	precision form \eqref{eq:mv-lik-diag}; its quadratic form can be taken
	over all $k$ components because the missing rows and columns of $\mP_i$
	vanish),
	\begin{align}
		p(\mR_\ell \mid \vmu_\ell, \mLam)
		 & = \prod_{i \in I_\ell}(2\pi)^{-k_i/2}
		\left|[\mA_i\mLam\mA_i]_{O_iO_i}\right|^{1/2}
		\exp\!\left(-\frac{1}{2}(\vr_i - \vmu_\ell)^\top\mP_i(\vr_i - \vmu_\ell)\right), \\
		p(\vmu_\ell \mid \mLamm)
		 & = (2\pi)^{-k/2}|\mLamm|^{1/2}
		\exp\!\left(-\frac{1}{2}\vmu_\ell^\top\mLamm\vmu_\ell\right).
	\end{align}
	Expanding $(\vr_i - \vmu_\ell)^\top\mP_i(\vr_i - \vmu_\ell)
		= \vr_i^\top\mP_i\vr_i - 2\vmu_\ell^\top\mP_i\vr_i
		+ \vmu_\ell^\top\mP_i\vmu_\ell$, summing over $i \in I_\ell$, and using
	\eqref{eq:mv-precision-sum}--\eqref{eq:mv-suffstats-B}, the joint exponent
	becomes
	\begin{equation}
		-\frac{1}{2}\sum_{i\in I_\ell}\vr_i^\top\mP_i\vr_i
		+ \vmu_\ell^\top\vB_\ell
		- \frac{1}{2}\vmu_\ell^\top\left(\mLamm + \mP_\ell\right)\vmu_\ell.
	\end{equation}
	With $V_\ell^{-1} = \mLamm + \mP_\ell$ and
	$\hat\vmu_\ell = V_\ell\vB_\ell$, completing the square in $\vmu_\ell$
	splits the last two terms as
	\begin{equation}
		-\frac{1}{2}(\vmu_\ell - \hat\vmu_\ell)^\top V_\ell^{-1}
		(\vmu_\ell - \hat\vmu_\ell) + \frac{1}{2}\vB_\ell^\top V_\ell\vB_\ell,
	\end{equation}
	and the Gaussian integral over $\R^k$ evaluates to
	\begin{align}
		\int_{\R^k}\exp\!\left(-\frac{1}{2}(\vmu_\ell - \hat\vmu_\ell)^\top
		V_\ell^{-1}(\vmu_\ell - \hat\vmu_\ell)\right) d\vmu_\ell
		 & = (2\pi)^{k/2}|V_\ell|^{1/2}                        \\
		 & = (2\pi)^{k/2}\bigl|\mLamm + \mP_\ell\bigr|^{-1/2}.
	\end{align}
	Collecting all prefactors,
	\begin{align}
		p(\mR_\ell \mid T, \mLam)
		 & = \prod_{i\in I_\ell}\left[(2\pi)^{-k_i/2}
		\left|[\mA_i\mLam\mA_i]_{O_iO_i}\right|^{1/2}\right]
		\cdot (2\pi)^{-k/2}|\mLamm|^{1/2}                                           \\
		 & \quad \cdot (2\pi)^{k/2}\,\bigl|\mLamm + \mP_\ell\bigr|^{-1/2}           \\
		 & \quad \cdot \exp\!\left(-\frac{1}{2}\sum_{i\in I_\ell}
		\vr_i^\top\mP_i\vr_i + \frac{1}{2}\vB_\ell^\top V_\ell\vB_\ell\right).
	\end{align}
	The $(2\pi)$ factors combine to $(2\pi)^{-\sum_{i\in I_\ell}k_i/2}$;
	taking the logarithm yields
	\eqref{eq:mv-marglik-a}--\eqref{eq:mv-marglik-c}.
\end{proof}

\paragraph{Reduced form for tree-structure moves.}
Equation \eqref{eq:mv-marglik-a} and the term
$-\frac{1}{2}\sum_{i\in I_\ell}\vr_i^\top\mP_i\vr_i$ are \emph{additive} in
the index set $I_\ell$: their sum across the leaves of any partition of a
fixed index set is the same. In particular, when comparing a leaf $\eta$
against its split $\eta \to (\eta_L, \eta_R)$ with $I_{\eta_L} \sqcup
	I_{\eta_R} = I_\eta$, those terms cancel exactly, and only the
log-determinant and quadratic-form pieces remain:
\begin{multline}
	\label{eq:mv-reduced-marglik}
	\log p(\mR_\ell \mid T, \mLam) =
	\underbrace{\frac{1}{2}\log|\mLamm|
		- \frac{1}{2}\log\bigl|\mLamm + \mP_\ell\bigr|}_{\text{log-det term}}
	\\
	+ \underbrace{\frac{1}{2}\,\vB_\ell^\top\bigl(\mLamm + \mP_\ell
		\bigr)^{-1}\vB_\ell}_{\text{quadratic term}}
	+ (\text{cancels in splits / prunes}).
\end{multline}
In bartz these two terms are assembled in \texttt{mcmcstep/\_step.py} by
\texttt{\_precompute\_likelihood\_terms\_mv\_het} for general $\va_i$, and by
\texttt{\_precompute\_likelihood\_terms\_mv} when $\va_i \equiv \Ind_k$ or
the scales are scalar per datapoint, so that $\mP_\ell \propto \mLam$. Note
that the
data-only statistics $\mN_\ell$ and $S_\ell$ in
\eqref{eq:mv-suffstats-N}--\eqref{eq:mv-suffstats-S} are additive across
partitions, so the relevant cancellations in the GROW/PRUNE comparisons
involve only the six quantities
$(\mN_{\eta_L}, \mN_{\eta_R}, \mN_\eta, S_{\eta_L}, S_{\eta_R}, S_\eta)$ with
$\mN_\eta = \mN_{\eta_L} + \mN_{\eta_R}$ and
$S_\eta = S_{\eta_L} + S_{\eta_R}$ (and likewise for the derived
$\vB_\eta = \vB_{\eta_L} + \vB_{\eta_R}$).

%-----------------------------------------------------------------------------
\section{Acceptance ratios for GROW and PRUNE}
\label{sec:mv-lkratios}

The MCMC sampler updates each tree by proposing a structural move on $T$
\emph{after marginalising over $M$}: the marginal posterior of $T$
conditional on $\mLam$ and the partial residuals is the target. This is the
standard ``integrate-the-leaves'' trick; sampling $M$ is then deferred to
the post-acceptance step using Proposition~\ref{prop:mv-leaf}.

\subsection{The GROW move}
\label{sec:mv-grow}

Pick a leaf $\eta \in L(T)$ and replace it by an internal node with a fresh
splitting rule and two leaf children $\eta_L, \eta_R$. Call the new tree $T'$.
The Metropolis--Hastings acceptance ratio for this move, with the leaves of
both $T$ and $T'$ marginalised, is
\begin{equation}
	\label{eq:mv-mh}
	\mathrm{MH}_{\mathrm{grow}}
	= \underbrace{\frac{q(T \mid T')}{q(T' \mid T)}}_{\text{transition}}
	\cdot \underbrace{\frac{p(T')}{p(T)}}_{\text{prior}}
	\cdot \underbrace{\frac{p(\vr_{1:n} \mid T', \mLam)}{p(\vr_{1:n} \mid T, \mLam)}}_{\text{marginal density}}.
\end{equation}

\paragraph{Marginal density ratio.}
Only the leaf $\eta$ (in $T$) and the leaves $\eta_L, \eta_R$ (in $T'$)
differ between the two collapsed densities; every other leaf has identical
sufficient statistics in $T$ and $T'$. Using the reduced form
\eqref{eq:mv-reduced-marglik},
\begin{align}
	\log\frac{p(\vr_{1:n} \mid T', \mLam)}{p(\vr_{1:n} \mid T, \mLam)}
	 & = \frac{1}{2}\Bigl[
		                \log|\mLamm|
		                + \log\bigl|\mLamm + \mP_\eta\bigr|
		                \label{eq:mv-lkratio-grow-a}                                        \\
		                & \qquad\;\;
		                - \log\bigl|\mLamm + \mP_{\eta_L}\bigr|
		                - \log\bigl|\mLamm + \mP_{\eta_R}\bigr|
		                \Bigr]
	\label{eq:mv-lkratio-grow-b}                                                                \\
	 & \quad + \frac{1}{2}\Bigl[
		                      \vB_{\eta_L}^\top\bigl(\mLamm + \mP_{\eta_L}\bigr)^{-1}\vB_{\eta_L}
		                      + \vB_{\eta_R}^\top\bigl(\mLamm + \mP_{\eta_R}\bigr)^{-1}\vB_{\eta_R}
		                      \label{eq:mv-lkratio-grow-c}                                        \\
		                      & \qquad\;\;
		                      - \vB_\eta^\top\bigl(\mLamm + \mP_\eta\bigr)^{-1}\vB_\eta
		                      \Bigr],
	\label{eq:mv-lkratio-grow}
\end{align}
with $\mN_\eta = \mN_{\eta_L} + \mN_{\eta_R}$, $S_\eta = S_{\eta_L} + S_{\eta_R}$,
and hence $\mP_\eta = \mP_{\eta_L} + \mP_{\eta_R}$ and
$\vB_\eta = \vB_{\eta_L} + \vB_{\eta_R}$ by the additivity of
\eqref{eq:mv-suffstats-N}--\eqref{eq:mv-suffstats-S}. The four
log-determinants and three quadratic forms can be computed stably from
Cholesky factors of $\mLamm + \mP_*$. This is what bartz does:
\texttt{\_precompute\_likelihood\_terms\_mv\_het} gathers the factors and
their log-determinants at $\eta_L$, $\eta_R$, $\eta$, and
\texttt{\_compute\_likelihood\_ratio\_mv\_het} evaluates the quadratic forms
by triangular solves against $\vB_*$. When $\va_i \equiv \Ind_k$ the
statistics $\mN_*$ and $S_*$ collapse to scaled counts and sums, and
\texttt{\_precompute\_likelihood\_terms\_mv} /
\texttt{\_compute\_likelihood\_ratio\_mv} instead precompute the sandwich
$\mLam(\mLamm + \mP_*)^{-1}\mLam$ and contract it with the plain
residual sums.

\paragraph{Prior ratio.}
Let $d_\eta$ be the depth of $\eta$, so that $\eta_L, \eta_R$ are at depth
$d_\eta + 1$. The only nodes whose prior status changes between $T$ and
$T'$ are $\eta$ (leaf $\to$ internal) and $\eta_L, \eta_R$ (introduced as
leaves). Writing $p_d \coloneqq p_{\mathrm{NT}}(d)$ for brevity,
\begin{equation}
	\label{eq:mv-prior-ratio-grow}
	\frac{p(T')}{p(T)}
	= \underbrace{\frac{p_{d_\eta}}{1 - p_{d_\eta}}}_{\eta \text{ now internal}}
	\cdot \underbrace{(1 - p_{d_\eta + 1})^{2}}_{\eta_L,\eta_R \text{ are leaves}}
	\cdot \underbrace{\frac{1}{n_{\mathrm{var}}(\eta)} \cdot \frac{1}{n_{\mathrm{cut}}(\eta, v)}}_{\text{rule prior at } \eta},
\end{equation}
where $n_{\mathrm{var}}(\eta)$ is the number of variables that admit at least
one valid cutpoint at $\eta$ and $n_{\mathrm{cut}}(\eta, v)$ is the number of
valid cutpoints for the chosen variable $v$. (If a child $\eta_L$ or
$\eta_R$ would have no further admissible split, its $1 - p_{d_\eta+1}$
factor is replaced by $1$, since the prior assigns it leaf status with
probability $1$. In the bartz code, this is encoded by the children
entries of \texttt{Moves.lrt\_growable}, used in \texttt{complete\_ratio}.)

\paragraph{Transition ratio.}
The forward kernel at $T$ chooses GROW with probability $1/2$ (or $1$ if
PRUNE is impossible, i.e., $T$ is the root), then picks $\eta$ uniformly from
the growable leaves, then a variable uniformly from the
$n_{\mathrm{var}}(\eta)$ admissible ones, then a cutpoint uniformly from the
$n_{\mathrm{cut}}(\eta, v)$ admissible ones:
\[
	q(T' \mid T) = \frac{1}{2(\text{or }1)} \cdot
	\frac{1}{|G(T)|} \cdot
	\frac{1}{n_{\mathrm{var}}(\eta)} \cdot
	\frac{1}{n_{\mathrm{cut}}(\eta, v)},
\]
where $G(T)$ is the set of growable leaves. The reverse kernel at $T'$
chooses PRUNE with probability $1/2$ (or $1$ if $T'$ has no growable leaves
left), then picks the prunable internal node uniformly from the
$|P(T')|$ candidates ($P(T')$ being the set of internal nodes whose two
children are both leaves):
\[
	q(T \mid T') = \frac{1}{2(\text{or }1)} \cdot \frac{1}{|P(T')|}.
\]
Hence
\begin{equation}
	\label{eq:mv-trans-ratio-grow}
	\frac{q(T \mid T')}{q(T' \mid T)}
	= \frac{P_{\mathrm{prune}}(T')}{P_{\mathrm{grow}}(T)} \cdot
	\frac{|G(T)|}{|P(T')|} \cdot
	n_{\mathrm{var}}(\eta) \cdot n_{\mathrm{cut}}(\eta, v),
\end{equation}
with $P_{\mathrm{grow}}, P_{\mathrm{prune}} \in \{\frac{1}{2}, 1\}$ as just
described.

\paragraph{Combined.}
Multiplying \eqref{eq:mv-prior-ratio-grow} and
\eqref{eq:mv-trans-ratio-grow}, the factors $1/n_{\mathrm{var}}(\eta)$ and
$1/n_{\mathrm{cut}}(\eta, v)$ cancel, leaving the standard CGM form:
\begin{equation}
	\label{eq:mv-nonlk-grow}
	\boxed{
		\frac{q(T \mid T')\,p(T')}{q(T' \mid T)\,p(T)}
		= \frac{P_{\mathrm{prune}}(T')}{P_{\mathrm{grow}}(T)} \cdot
		\frac{|G(T)|}{|P(T')|} \cdot
		\frac{p_{d_\eta}}{1 - p_{d_\eta}} \cdot
		(1 - p_{d_\eta+1})^{2}.
	}
\end{equation}
This non-likelihood factor depends only on the tree topology and the CGM
prior, not on the leaf or noise model. It is precisely
\texttt{Moves.log\_trans\_prior\_ratio}: \texttt{compute\_partial\_ratio} in
\texttt{mcmcstep/\_moves.py} computes
$|G(T)| / (P_{\mathrm{grow}}(T) |P(T')|)$ (the variable / cutpoint
cancellation is the comment ``the transition and prior ratios also contain
factors num\_available\_split * num\_available\_var * s[var], but they cancel
out''), and \texttt{complete\_ratio} in \texttt{mcmcstep/\_step.py}
multiplies in $P_{\mathrm{prune}}(T')$ and the two prior factors.

\subsection{The PRUNE move}
\label{sec:mv-prune}

A PRUNE move acts as the inverse of a GROW: it picks an internal node
$\eta \in P(T)$ whose two children are leaves, deletes the children, and
turns $\eta$ into a leaf. Because the two moves form a reversible pair, the
acceptance ratio of a PRUNE from $T$ to $T'$ is simply the reciprocal of the
GROW ratio for the move $T' \to T$:
\begin{equation}
	\label{eq:mv-mh-prune}
	\mathrm{MH}_{\mathrm{prune}}
	= \mathrm{MH}_{\mathrm{grow}}^{\,-1}(T' \to T).
\end{equation}
In particular, the log marginal-density ratio is the negative of
\eqref{eq:mv-lkratio-grow}, with the sufficient statistics
$(\mN_{\eta_L}, \mN_{\eta_R}, S_{\eta_L}, S_{\eta_R})$ read off the
to-be-deleted children, and the non-likelihood factor
\eqref{eq:mv-nonlk-grow} inverted for the inverse move. Both \texttt{bartz} and \texttt{stochtree}
exploit this symmetry:
\begin{itemize}
	\item \texttt{bartz} evaluates the GROW ratio for both moves and negates
	      the log-ratio when the proposal is PRUNE
	      (\texttt{accept\_move\_and\_sample\_leaves} in
	      \texttt{mcmcstep/\_step.py}: \texttt{log\_ratio = jnp.where(pt.move.grow,
		      log\_ratio, -log\_ratio)}).
	\item \texttt{stochtree} compares
	      \texttt{NoSplitLogMarginalLikelihood} against
	      \texttt{SplitLogMarginalLikelihood} with sign reversed for prune
	      steps inside \texttt{sampler.cpp}.
\end{itemize}

\subsection{Sampling the leaves after acceptance}
\label{sec:mv-sample-leaves}

Whenever the structure update is finalised (whether the proposal was
accepted or rejected), the leaves of $T$ are resampled from
\eqref{eq:mv-leaf-fc}. In the sum-of-trees setting, the residuals of every
other tree are then refreshed by subtracting the new
$g(\cdot;\, T_t, M_t)$ from the running sum, and the cycle moves to the
next tree. After all $m$ trees have been updated, $\mLam$ is drawn from its
full conditional, derived next.

%-----------------------------------------------------------------------------
\section{Updating the error precision}
\label{sec:mv-Sigma}

After all $m$ trees have been updated, the global precision $\mLam$ is drawn
from its full conditional. Let
$\bm{\veps}_i \coloneqq \vy_i - f(\vx_i)$ be the residual of the full forest,
and let
\begin{equation}
	\label{eq:mv-eps-tilde}
	\tilde{\bm{\veps}}_i \coloneqq \mA_i \bm{\veps}_i
	= \va_i \odot \bm{\veps}_i
\end{equation}
be its rescaling, which vanishes on the missing components. By
\eqref{eq:mv-lik-diag},
$\tilde{\bm{\veps}}_{i,O_i} \mid \mLam \sim \N_{k_i}(\mathbf{0}, \mLam_{O_iO_i}^{-1})$
independently across $i$. Write
\begin{equation}
	\label{eq:mv-Sigma-rate}
	\Psi' \coloneqq \Psi + \sum_{i=1}^n \tilde{\bm{\veps}}_i \tilde{\bm{\veps}}_i^\top
	= \Psi + \sum_{i=1}^n \tilde\mA_i \odot (\bm{\veps}_i \bm{\veps}_i^\top)
\end{equation}
for the posterior rate matrix.

\begin{prop}[Error precision full conditional]
	\label{prop:mv-Sigma}
	Let $\Psi'$ be the posterior rate matrix \eqref{eq:mv-Sigma-rate}.
	\begin{enumerate}[label=(\roman*)]
		\item Without missing values ($O_i = \{1,\dots,k\}$ for all $i$),
		      \begin{equation}
			      \label{eq:mv-Sigma-fc}
			      \mLam \mid \cdots \sim \W\left(\nu + n, \Psi'\right).
		      \end{equation}
		\item If $\mLam = \diag(\lambda_1, \dots, \lambda_k)$ is diagonal, with
		      independent shape--rate gamma priors
		      $\lambda_j \sim \Gamma(\nu_j/2, \psi_j/2)$ in place of
		      \eqref{eq:mv-iw-prior}, the components stay independent a
		      posteriori,
		      \begin{equation}
			      \label{eq:mv-Sigma-fc-diag}
			      \lambda_j \mid \cdots \sim
			      \Gamma\left(\frac{\nu_j + n_j}{2},
			      \frac{\psi_j + \sum_{i=1}^n a_{i,j}^2 \veps_{i,j}^2}{2}\right),
			      \qquad
			      n_j \coloneqq |\{i : j \in O_i\}|,
		      \end{equation}
		      whatever the pattern of missing values.
	\end{enumerate}
\end{prop}

\begin{proof}
	The joint log density of $\mLam$ and the observed rescaled residuals,
	treating $\tilde{\bm{\veps}}_{1:n}$ as observed, is
	\begin{align}
		\log p(\tilde{\bm{\veps}}_{1:n} \mid \mLam) + \log p(\mLam)
		 & = \sum_{i=1}^n \left[
		\frac{1}{2}\log\left|\mLam_{O_iO_i}\right|
		- \frac{1}{2}\tilde{\bm{\veps}}_i^\top \mLam \tilde{\bm{\veps}}_i
		\right]        \\
		 & \quad
		+ \frac{\nu - k - 1}{2}\log|\mLam|
		- \frac{1}{2}\tr(\Psi \mLam) + \text{const},
	\end{align}
	where the quadratic form is taken over all $k$ components because
	$\tilde{\bm{\veps}}_i$ vanishes off $O_i$. Writing
	$\tilde{\bm{\veps}}_i^\top \mLam \tilde{\bm{\veps}}_i =
		\tr(\tilde{\bm{\veps}}_i\tilde{\bm{\veps}}_i^\top \mLam)$ and collecting
	the trace terms into $\Psi'$, this is
	\begin{equation}
		\frac{\nu - k - 1}{2}\log|\mLam|
		+ \frac{1}{2}\sum_{i=1}^n \log\left|\mLam_{O_iO_i}\right|
		- \frac{1}{2}\tr(\Psi' \mLam) + \text{const}.
	\end{equation}
	For (i), $|\mLam_{O_iO_i}| = |\mLam|$ for every $i$, so the exponent of
	$|\mLam|$ is $(\nu + n - k - 1)/2$, the kernel of $\W(\nu + n, \Psi')$.
	For (ii), the prior term is replaced by
	$\sum_j [(\nu_j/2 - 1)\log\lambda_j - \psi_j\lambda_j/2]$,
	$|\mLam_{O_iO_i}| = \prod_{j \in O_i}\lambda_j$, and
	$\tr(\Psi'\mLam) = \sum_j \Psi'_{jj}\lambda_j$ with
	$\Psi'_{jj} = \psi_j + \sum_i a_{i,j}^2\veps_{i,j}^2$, so the density
	factorises over $j$ as
	$\prod_j \lambda_j^{(\nu_j + n_j)/2 - 1}\exp(-\Psi'_{jj}\lambda_j/2)$.
\end{proof}

The proof also shows why the diagonal restriction of \S\ref{sec:mv-model}
is needed: with missing values, $\prod_i |\mLam_{O_iO_i}|^{1/2}$ is not a
power of $|\mLam|$, so even the likelihood \eqref{eq:mv-lik-diag} would not
be conjugate to the Wishart prior for a dense $\mLam$, and the actual
marginal likelihood \eqref{eq:mv-lik-obs} involves the Schur complements
$[(\mLam^{-1})_{O_iO_i}]^{-1}$ on top of that. For $k = 1$, (ii) is the
univariate update of \S\ref{sec:sigma-update} with $n$ replaced by the
number of observed outcomes.

Note that the heteroskedasticity is handled \emph{exactly} by the residual
rescaling $\bm{\veps}_i \mapsto \va_i \odot \bm{\veps}_i$: no auxiliary
sampler or block update is required.

\paragraph{Sampling.} A draw from $\W(\nu', \Psi')$ is obtained directly,
e.g.\ by a Bartlett decomposition of the Cholesky factor of $(\Psi')^{-1}$;
this is what \texttt{\_sample\_wishart\_bartlett} and
\texttt{\_step\_error\_cov\_inv\_mv} do in bartz (\texttt{mcmcstep/\_step.py}),
the latter rescaling the residuals by \texttt{State.inv\_sdev\_scale}
$= \va_i$ before forming $\Psi' = \Psi + \tilde R^\top \tilde R$, with
$\tilde R$ the row-rescaled residual matrix. The diagonal update
\eqref{eq:mv-Sigma-fc-diag} is \texttt{\_step\_error\_cov\_inv\_diag}, with
the prior represented by \texttt{DiagWishart} ($\nu_j = \nu$ common to all
components, $\psi_j = \Psi_{jj}$); \texttt{init} requires it when the
\texttt{missing} mask is per component.

%-----------------------------------------------------------------------------
\section{Implementation pointers}
\label{sec:impl-full}

Bartz implements this chapter as the path selected by
\texttt{mcmcstep.init} when \texttt{error\_scale} is a $k \times n$ array of
per-component scales $1/\va_i$ (a $k \times n$ \texttt{missing} mask sets
the corresponding $a_{i,j} = 0$ and, as in \S\ref{sec:mv-model}, requires a
diagonal $\mLam$, i.e.\ a \texttt{DiagWishart} prior). The relevant files
are all in
\texttt{src/bartz/mcmcstep}.
\begin{description}
	\item[\texttt{\_state.py}] The \texttt{State} and \texttt{Forest} classes
	      (state and priors):
	      \begin{itemize}
		      \item \texttt{State.error\_cov\_inv} is $\mLam$ with its
		            \texttt{Wishart} prior (\texttt{nu} $= \nu$, \texttt{rate}
		            $= \Psi$);
		      \item \texttt{Forest.leaf\_prior\_cov\_inv} is $\mLamm$;
		      \item \texttt{State.inv\_sdev\_scale} is $\va_i$ and
		            \texttt{State.prec\_scale} is $\tilde\mA_i$ (both stored
		            up to a power-of-two unit);
		      \item \texttt{Forest.prec\_tree} holds $\mN_\ell$.
	      \end{itemize}
	\item[\texttt{\_moves.py}] GROW/PRUNE proposals and
	      \texttt{compute\_partial\_ratio}. These implement
	      \S\ref{sec:mv-grow}--\ref{sec:mv-prune} and do not depend on the
	      leaf model.
	\item[\texttt{\_step.py}] The leaf and noise model:
	      \begin{itemize}
		      \item \texttt{complete\_ratio} finishes \eqref{eq:mv-nonlk-grow};
		      \item \texttt{sum\_resid} accumulates $S_\ell$ and
		            \texttt{compute\_B} forms $\vB_\ell$;
		      \item \texttt{\_precompute\_leaf\_terms\_mv\_het} computes
		            \eqref{eq:mv-leaf-fc};
		      \item \texttt{\_precompute\_likelihood\_terms\_mv\_het} and
		            \texttt{\_compute\_likelihood\_ratio\_mv\_het} compute
		            \eqref{eq:mv-lkratio-grow};
		      \item \texttt{\_step\_error\_cov\_inv\_mv} and
		            \texttt{\_sample\_wishart\_bartlett} draw from
		            \eqref{eq:mv-Sigma-fc};
		      \item \texttt{\_step\_error\_cov\_inv\_diag} draws from
		            \eqref{eq:mv-Sigma-fc-diag}.
	      \end{itemize}
	      When $\mP_\ell \propto \mLam$ (homoskedastic, or scalar scale per
	      datapoint) the same steps are dispatched to the cheaper
	      \texttt{\_mv} variants of these functions, with $n_\ell$ from
	      \texttt{Forest.count\_tree}.
\end{description}
The Dirichlet prior on the split variable (\texttt{Forest.log\_s}) and
probit outcomes are not covered here.
The tree-topology parts (tree prior, proposal kernel, non-likelihood ratio
\eqref{eq:mv-nonlk-grow}) are shared verbatim with the univariate code of
Chapter~\ref{ch:vanilla}.
